\section{Convex combinations}
A linear combination allows arbitrary coefficients. An affine combination
requires that their sum be $1$. A convex combination requires both sum $1$
and \emph{nonnegative} coefficients. Thus
$\operatorname{conv}S\subseteq\operatorname{aff}S\subseteq\Span S$.
For an affinely independent list, its affine coordinates are unique: a
negative coordinate rules out convex membership, all positive coordinates
put the point in the relative interior of its simplex, and zero coordinates
locate it on a face. A convex hull consists of finite convex combinations;
it need not include limit points not obtained by those combinations.

\practice{8.3.P1}{Test a triangle in three-dimensional space}
Use the supplied vertices $v_1=(6,2,2)$, $v_2=(7,1,5)$, $v_3=(-2,4,-1)$.
To find affine coordinates solve
$p-v_1=a(v_2-v_1)+b(v_3-v_1)$, then use coefficients $(1-a-b,a,b)$.
For $p_1=(1,3,1)$ the solution is $a=1/3,b=2/3$; checking gives
$p_1=0v_1+(1/3)v_2+(2/3)v_3$. All coefficients are nonnegative, so
$p_1\in\operatorname{conv}S$.
For $p_2=(3,2,1)$ the first two coordinate equations give $a=1,b=1/2$,
but the third would require $3a-3b=-1$, whereas it equals $3/2$.
Thus $p_2$ is not even in the affine hull and cannot be in the convex hull.

\practice{8.3.P2}{The hull of the positive reciprocal curve}
\goal Describe the hull of $S=\{(x,1/x):x>0\}$ exactly.
No convex combination has $x\le0$. Moreover convexity of $1/x$ on positive
inputs gives $\sum_i\lambda_i/x_i\ge1/(\sum_i\lambda_ix_i)$.
For example this follows directly from Cauchy--Schwarz applied to
$\sqrt{\lambda_ix_i}$ and $\sqrt{\lambda_i/x_i}$.
Hence the hull lies in $\{(x,y):x>0,y\ge1/x\}$.
Conversely, for a point $(x,y)$ in that set draw the line of slope $-1$
through it, $u+v=x+y$. Its intersections with $uv=1$ have positive
$x$-coordinates given by roots of $u^2-(x+y)u+1=0$.
The polynomial at $u=x$ is $1-xy\le0$, so $x$ lies between the two
roots. Thus $(x,y)$ lies on the line segment joining these two curve
points, or is itself the single tangency point. It is a convex combination
of them. Therefore the hull is exactly the region on and above the curve.

\exercise{8.3.1}{A triangle with a missing upper edge}
Every convex combination of vertical-segment points and $(2,0)$ can be
written $(1-\lambda)(0,a)+\lambda(2,0)$ with
$0\le\lambda\le1$ and $0\le a<1$. (Combine all the vertical points
into their weighted average first.) Thus the hull is
\[
 \{(x,y):0\le x<2,\ 0\le y<1-x/2\}\ \cup\ \{(2,0)\}.
\]
Conversely any point in the first set uses $\lambda=x/2$ and
$a=y/(1-x/2)$, which satisfy the required inequalities.
The lower edge and right vertex are included, but the sloping upper edge
and the point $(0,1)$ are excluded. Replacing the strict inequality by a
weak one would add points not actually in the convex hull.

\exercise{8.3.2}{Three convex hulls of graphs}
\step{(a) The reciprocal graph starting at $x=1/2$.}
The exact hull is
\[
 \{(1/2,2)\}\ \cup\ \{(x,y):x>1/2,\ 1/x\le y<2\}.
\]
Every combination has $x\ge1/2$, lies on or above the reciprocal curve by
convexity, and has $y\le2$. Equality $y=2$ forces every positive-weight
input to be the endpoint $(1/2,2)$, so it occurs only there.
For the converse fix $x>1/2$ and $1/x\le y<2$. Put
$b=(2x-1)/(2-y)$. The lower bound on $y$ gives $b\ge x$.
The chord joining $(1/2,2)$ and $(b,1/b)$ has height
$2-(2x-1)/b=y$ at horizontal coordinate $x$.
Because $1/2<x\le b$, the point lies on that segment and is in the hull.
\step{(b) The sine graph on the whole real line.}
The hull is the closed strip $\{(x,y):-1\le y\le1\}$.
The range of sine bounds every combination within the strip.
Conversely its maximum points $(\pi/2+2j\pi,1)$ extend in both horizontal
directions. For any fixed real $x$, two such points bracket $x$, so their
segment contains $(x,1)$. The minimum points similarly give $(x,-1)$.
The segment between these two hull points contains every $(x,y)$ in the strip.
\step{(c) The square-root graph.}
The hull is
\[
 \{(0,0)\}\ \cup\ \{(x,y):x>0,\ 0<y\le\sqrt{x}\}.
\]
Concavity of square root bounds the height of a convex combination by
$\sqrt{x}$; nonnegativity is automatic. If $x>0$, some positive-weight
input has positive $x$ and therefore positive height, so height cannot be zero.
Conversely, for $x>0$ and $0<y\le\sqrt{x}$, take
$a=(x/y)^2$ and $\lambda=y^2/x\in(0,1]$.
Then $\lambda(a,\sqrt a)+(1-\lambda)(0,0)=(x,y)$.
This proves the formula and explains why the positive horizontal axis is not included.

\exercise{8.3.3}{Use the unique coordinates from Section 8.1, Exercise 5}
The three orthogonal generators are
$b_1=(2,1,1)$, $b_2=(1,0,-2)$, $b_3=(2,-5,1)$.
Their squared lengths are $6,5,30$, so the unique linear coefficients of
$p$ are $c_j=(p\cdot b_j)/(b_j\cdot b_j)$.
For the three supplied points the results are
\[
 \begin{array}{c|c|c}
 p& (c_1,c_2,c_3)&\text{sum}\\\hline
 (3,8,4)&(3,-1,-1)&1\\
 (6,-3,3)&(2,0,1)&3\\
 (0,-1,-5)&(-1,2,0)&1
 \end{array}
\]
The first and third have negative coordinates; the second does not even
have coefficient sum one. Since the linear representation in a basis is
unique, no alternative nonnegative sum-one representation exists.
\answer{None of the three points is in $\operatorname{conv}S$.}

\exercise{8.3.4}{A second list with the same orthogonal generators}
The same dot-product calculation gives
\[
 \begin{array}{c|c|c}
 p&(c_1,c_2,c_3)&\text{sum}\\\hline
 (0,-19,-5)&(-4,2,3)&1\\
 (3/2,-13/10,-1/2)&(1/5,1/2,3/10)&1\\
 (5,-4,0)&(1,1,1)&3
 \end{array}
\]
Only the middle row has all coefficients nonnegative and sum one.
Thus $p_2\in\operatorname{conv}S$, while $p_1,p_3$ are not.

\exercise{8.3.5}{Convex coordinates in a tetrahedron}
Append a final coordinate $1$ to each of the four supplied $v$-vectors.
Solving the resulting $4\times4$ systems for the two points gives
\[
 [p_1]_S=(-1/6,1/3,2/3,1/6),\qquad
 [p_2]_S=(1/3,1/3,1/6,1/6).
\]
Here coordinates refer to affine/barycentric coordinates. The appended
matrix is invertible, so these representations are unique.
For instance $(1/3)v_1+(1/3)v_2+(1/6)v_3+(1/6)v_4=(0,-2,2)=p_2$.
The first representation has a negative coefficient and therefore is not
convex; the second has all coefficients positive and sum one.
\answer{$p_1$ is outside the convex hull; $p_2$ is in it.}

\exercise{8.3.6}{Distinguish the linear, affine, and convex hulls}
The supplied orthogonal vectors are
$v_1=(2,0,-1,2)$, $v_2=(0,-2,2,1)$, $v_3=(-2,1,0,2)$;
each has squared length $9$.
For a proposed $p$, first compute $c_j=(p\cdot v_j)/9$ and check whether
$\sum c_jv_j=p$. Only after this check do sum-one and sign tests apply.
\begin{enumerate}[label=(\alph*)]
\item For $p_1=(-1,2,-3/2,5/2)$, the coefficients are $(1/2,-1/2,1)$
and reconstruct $p_1$. They sum to one but include a negative coefficient.
Thus $p_1$ is in the span and affine hull, but not the convex hull.
\item For $p_2=(-1/2,0,1/4,7/4)$, the coefficients are $(1/4,1/4,1/2)$.
They reconstruct $p_2$, sum to one, and are nonnegative. Thus $p_2$ is in
all three hulls.
\item For $p_3=(6,-4,1,-1)$, the coefficients are $(1,1,-2)$ and
reconstruct $p_3$, but their sum is zero. Hence $p_3$ is in the span only,
not the affine or convex hull.
\item For $p_4=(-1,-2,0,4)$, the dot-product coefficients are
$(2/3,8/9,8/9)$, but the reconstruction error is
$p_4-\sum c_jv_j=(-5/9,-10/9,-10/9,0)\ne0$.
Thus $p_4$ is not in the span and consequently in neither smaller hull.
\end{enumerate}

\exercise{8.3.7}{Barycentric coordinates locate four planar points}
\step{(a) Solve the affine-coordinate systems.}
With $T=\{(-1,0),(2,3),(4,1)\}$, solve
\[
 \begin{pmatrix}-1&2&4\\0&3&1\\1&1&1\end{pmatrix}c
 =\begin{pmatrix}x\\y\\1\end{pmatrix}.
\]
The four answers, in the listed order, are
\[
 (1/3,1/6,1/2),\quad(0,1/2,1/2),\quad
 (1/2,-1/4,3/4),\quad(1/2,3/4,-1/4).
\]
\step{(b) Use signs and zero entries.}
$p_1$ has all positive entries and is inside the triangle.
$p_2$ has one zero and the other entries positive, so it is on the edge
joining vertices two and three. $p_3,p_4$ have negative entries and are outside.

\exercise{8.3.8}{Another barycentric-coordinate table}
For $T=\{(2,0),(0,5),(-1,1)\}$, use the appended-coordinate system
\[
 \begin{pmatrix}2&0&-1\\0&5&1\\1&1&1\end{pmatrix}c=(x,y,1)^T.
\]
For $p_1,p_2,p_3,p_4$ its solutions are respectively
\[
 (12/13,3/13,-2/13),\quad(8/13,2/13,3/13),\quad
 (2/3,0,1/3),\quad(9/13,-1/13,5/13).
\]
Thus $p_1$ is outside, $p_2$ inside, $p_3$ on the edge joining vertices
one and three, and $p_4$ outside. The $p_3$ supplied in the exercise is
$(1,1/3)$, not $(1,1)$; preserving this fraction is essential.

\exercise{8.3.9}{Read positions from tetrahedral coordinates}
A point is inside precisely when all four barycentric coordinates are
positive. Nonnegative coordinates with a zero put it on a surface; two
zeros put it on an edge. A negative coordinate puts it outside.
Therefore $p_1=(2,0,0,-1)$ is outside;
$p_2=(0,1/2,1/4,1/4)$ is on the face opposite vertex one, but not an edge;
$p_3=(1/2,0,3/2,-1)$ is outside;
$p_4=(1/3,1/4,1/4,1/6)$ is inside; and
$p_5=(1/3,0,2/3,0)$ is on the edge joining vertices one and three.

\exercise{8.3.10}{A second tetrahedral list}
Using the same criteria, $q_1=(1/8,1/4,1/8,1/2)$ is inside;
$q_2=(3/4,-1/4,0,1/2)$ is outside;
$q_3=(0,3/4,1/4,0)$ is on the edge joining vertices two and three;
$q_4=(0,-2,0,3)$ is outside; and
$q_5=(1/3,1/3,1/3,0)$ is on the face opposite vertex four, not an edge.

\exercise{8.3.11}{Three statements about convexity}
\begin{enumerate}[label=(\alph*)]
\item \textbf{False.} Sum one makes a combination affine, but negative
coefficients are still possible. Convexity also requires each coefficient nonnegative.
\item \textbf{False.} A singleton $S=\{v\}$ has $\operatorname{conv}S=S$.
More generally this equality holds for any convex $S$.
\item \textbf{False.} Two singleton sets are convex but their union
$\{(0,0),(2,0)\}$ misses the midpoint $(1,0)$.
\end{enumerate}

\exercise{8.3.12}{Three more statements about hulls}
\begin{enumerate}[label=(\alph*)]
\item \textbf{True.} Containing every segment between two of its points
is exactly the definition of convexity.
\item \textbf{True.} If two points belong to both convex sets, their
whole joining segment belongs to both, hence to the intersection. The empty
intersection is also convex under the usual vacuous definition.
\item \textbf{False.} Carath\'eodory's theorem says \emph{at most} six
points in $\R^5$, not necessarily six distinct available points.
A singleton $S$ is an immediate counterexample to the stated requirement.
\end{enumerate}

\exercise{8.3.13}{A linear image preserves convexity}
For $u=f(x),v=f(y)$ in $f(S)$ and $0\le t\le1$, linearity gives
$(1-t)u+tv=f((1-t)x+ty)$.
Convexity of $S$ puts the input $(1-t)x+ty$ in $S$, so this output is in
$f(S)$. Hence $f(S)$ is convex. Empty sets cause no exception.

\exercise{8.3.14}{A linear preimage preserves convexity}
For $x,y\in S$, the images $f(x),f(y)$ lie in $T$.
Their combination $(1-t)f(x)+tf(y)$ lies in $T$ for $0\le t\le1$ by
convexity. This equals $f((1-t)x+ty)$ by linearity, so the input
$(1-t)x+ty$ lies in $S$. Thus the whole segment is in $S$.

\exercise{8.3.15}{Eliminate a vertex in two different ways}
Direct coordinate addition checks both supplied identities:
$p=(1/3)v_1+(1/3)v_2+(1/6)v_3+(1/6)v_4=(2,1)$ and
$v_1-v_2+v_3-v_4=0$.
Add $s$ times the zero relation to the coefficients. They become
\[
 (1/3+s,\ 1/3-s,\ 1/6+s,\ 1/6-s).
\]
Their sum stays one and their combination stays $p$. For $s=1/6$ and
$s=-1/6$ all entries are nonnegative and one is zero, giving respectively
\[
 p=\tfrac12v_1+\tfrac16v_2+\tfrac13v_3,
 \qquad p=\tfrac16v_1+\tfrac12v_2+\tfrac13v_4.
\]
This is the coefficient-elimination step in the proof of Carath\'eodory's theorem.

\exercise{8.3.16}{A second coefficient-elimination example}
The printed instruction says ``Repeat Exercise 9,'' but the supplied
convex combination and zero relation clearly call for the procedure of
Exercise 15. We carry out that intended task explicitly.
The given coefficients are $(1,72,37,11)/121$, whose sum is one;
coordinate substitution gives $p=(1,2)$. The relation
$10v_1-6v_2+7v_3-11v_4=0$ also has coefficient sum zero.
Add $s(10,-6,7,-11)$ to the original coefficients.
For $s=1/121$, the last coefficient becomes zero and we obtain
\[
 p=\tfrac1{11}v_1+\tfrac6{11}v_2+\tfrac4{11}v_3.
\]
For $s=-1/1210$, the first becomes zero and the result is
\[
 p=\tfrac35v_2+\tfrac3{10}v_3+\tfrac1{10}v_4.
\]
Each list is nonnegative, sums to one, and reconstructs $p$.

\exercise{8.3.17}{A containing convex set contains the hull}
If $A\subseteq B$ and $B$ is convex, then any finite convex combination
of points in $A$ is a convex combination of points in $B$ and therefore
belongs to $B$. Every point of $\operatorname{conv}A$ is such a combination,
so $\operatorname{conv}A\subseteq B$.

\exercise{8.3.18}{Convex hull is monotone under inclusion}
If $A\subseteq B$, every convex combination of points in $A$ is also one
of points in $B$. Therefore $\operatorname{conv}A\subseteq\operatorname{conv}B$.
This also follows from the preceding exercise because $\operatorname{conv}B$
is convex and contains $A$.

\exercise{8.3.19}{Union of hulls need not fill the combined hull}
\step{(a) Prove the inclusion.}
Since $A\subseteq A\cup B$ and $B\subseteq A\cup B$, monotonicity gives
both $\operatorname{conv}A\subseteq\operatorname{conv}(A\cup B)$ and
the corresponding inclusion for $B$. Taking their union proves the claim.
\step{(b) Show strict inclusion can occur.}
Let $A=\{(0,0)\}$ and $B=\{(2,0)\}$.
The union of their hulls has just two points, whereas the combined hull is
the whole segment and includes $(1,0)$.

\exercise{8.3.20}{Intersection of hulls can be larger}
\step{(a) Apply monotonicity twice.}
Since $A\cap B\subseteq A$ and $A\cap B\subseteq B$, its hull is
contained in each hull, hence in their intersection.
\step{(b) Give an example with a new intersection point.}
Take $A=\{(-1,0),(1,0)\}$ and $B=\{(0,-1),(0,1)\}$.
The sets have empty intersection, so its hull is empty. But their two
segments intersect at zero, showing that the claimed inclusion can be strict.

\exercise{8.3.21}{The midpoint construction for a quadratic curve}
At $t=1/2$,
\[
 f_0=\tfrac12(p_0+p_1),\quad f_1=\tfrac12(p_1+p_2),\quad
 g=\tfrac12(f_0+f_1)=\tfrac14p_0+\tfrac12p_1+\tfrac14p_2.
\]
Thus first mark the midpoints of $p_0p_1$ and $p_1p_2$, join them,
and mark the midpoint of that new segment. The original diagram below
uses the same right-angle configuration as the source, with convenient
schematic coordinates $p_0=(0,0),p_1=(0,2),p_2=(4,2)$.
\begin{center}\setlength{\unitlength}{1mm}\begin{picture}(92,49)
\put(14,7){\polyline(0,0)(0,30)(60,30)}
\put(14,7){\color{gray}\polyline(0.000,15.000)(30.000,30.000)}
\put(14.000,7.000){\circle*{1.4}}
\put(5.000,5.000){$p_0$}
\put(14.000,37.000){\circle*{1.4}}
\put(5.000,39.000){$p_1$}
\put(74.000,37.000){\circle*{1.4}}
\put(76.000,39.000){$p_2$}
\put(14.000,22.000){\circle*{1.4}}
\put(3.000,21.000){$f_0$}
\put(44.000,37.000){\circle*{1.4}}
\put(46.000,39.000){$f_1$}
\put(29.000,29.500){\circle*{1.4}}
\put(30.000,22.500){$g$}
\end{picture}\end{center}

\exercise{8.3.22}{A three-quarter construction}
Now $f_0=(1/4)p_0+(3/4)p_1$ and $f_1=(1/4)p_1+(3/4)p_2$.
Place each point three quarters of the way along its directed segment.
Then place $g$ three quarters of the way from $f_0$ to $f_1$:
\[
 g=\tfrac1{16}p_0+\tfrac38p_1+\tfrac9{16}p_2.
\]
The same schematic coordinates give $f_0=(0,3/2)$, $f_1=(3,2)$,
and $g=(9/4,15/8)$, as shown.
\begin{center}\setlength{\unitlength}{1mm}\begin{picture}(92,49)
\put(14,7){\polyline(0,0)(0,30)(60,30)}
\put(14,7){\color{gray}\polyline(0.000,22.500)(45.000,30.000)}
\put(14.000,7.000){\circle*{1.4}}
\put(5.000,5.000){$p_0$}
\put(14.000,37.000){\circle*{1.4}}
\put(5.000,39.000){$p_1$}
\put(74.000,37.000){\circle*{1.4}}
\put(76.000,39.000){$p_2$}
\put(14.000,29.500){\circle*{1.4}}
\put(3.000,28.500){$f_0$}
\put(59.000,37.000){\circle*{1.4}}
\put(61.000,39.000){$f_1$}
\put(47.750,35.125){\circle*{1.4}}
\put(48.750,28.125){$g$}
\end{picture}\end{center}

\exercise{8.3.23}{Expand the quadratic B\'ezier formula}
Substitute both definitions into $g=(1-t)f_0+tf_1$ and distribute:
\[
 g(t)=(1-t)^2p_0+2t(1-t)p_1+t^2p_2.
\]
For $0\le t\le1$, the three coefficients are nonnegative and their sum
is $(1-t)^2+2t(1-t)+t^2=((1-t)+t)^2=1$.
Thus this formula is a convex combination of the control points and lies
in their convex hull for the entire parameter interval.

\exercise{8.3.24}{The cubic curve also stays inside its controls}
Apply the preceding formula to the two quadratic curves and collect terms:
\[
 h(t)=(1-t)^3p_0+3t(1-t)^2p_1+3t^2(1-t)p_2+t^3p_3.
\]
The four coefficients are nonnegative on $[0,1]$ and add to
$((1-t)+t)^3=1$ by the binomial theorem.
Hence every point of the cubic curve is a convex combination of the four
controls and lies in their convex hull. Equivalently, both quadratic
points already lie in that convex hull, which is itself convex.
